%!TEX root = /Users/markelikalderon/Documents/Git/perceptual_self-consciousness/perceptual_self-consciousness.tex
\chapter{Procession and Reversion} % (fold)
\label{cha:procession}

\section{Philoponus' Challenge} % (fold)
\label{sec:philoponus_challenge}

\begin{quote}
	Besides, it is absurd that the same sense should know that it sees. For it must be reflecting back on itself after having seen the colour that it gets to know that it sees. But if it reflects on itself it also has an activity which is separate, and what has a sperate activity has a separate substance, and on that account is eternal and incorporeal. So someone who says in an ambiguous way tha the rational soul is immortal will be shouting out plainly that the non-rational soul is immortal; which is absurd. The senses are not eternal, and for that reason do not reflect on themselves. And if they do not reflect on themselves they do not lay hold of their own activities. For a thing's reflecting on itself is nothing other than the laying hold of itself. (Philoponus)
\end{quote}

\begin{quotation}
	Prop. 15: \emph{All that is capable of reverting upon itself is incorporeal}
	For it is not in that nature of any body to revert upon itself. That which reverts upon anything is conjoined with that upon which it reverts: hence it is evident that every part of a body reverted upon itself must be conjoined with every other part—since self-reversion is precisely the case in which the reverted subject and that upon which it has reverted become identical. But this is impossible for a body, and universally for any divisible substance: for the whole of a divisible substance cannot be conjoined with the whole of itself, because of the separation of its parts, which occupy different positions in space. It is not in the nature, then, of any body to revert upon itself so that the whole is reverted upon the whole. Thus if there is anything which is capable of reverting upon itself, it is incorporeal and without parts. (Proclus, \emph{Elements of Theology}; \citealt[19]{Dodds:1963ul})
\end{quotation}


% section philoponus_challenge (end)


% Chapter procession (end)